% BEGIN REV09_MEMORIA_PROFUNDA
\subsection{Coordenada profunda y recuperación entera}
\label{corpus9:xii:memoria-profunda}

La coordenatización de seis coordenadas conserva el residuo expresado y el cociente
entero que determina su prolongación. Fijamos el representante entero de
\(L_0\) mostrado en la construcción regional:
\begin{equation}
 A_H=\begin{pmatrix}
 2&2&2&1&2&1\\
 2&1&2&2&1&1\\
 1&0&0&1&0&2\\
 0&0&1&1&1&0\\
 2&2&2&0&2&1\\
 2&1&2&1&0&0
 \end{pmatrix},\qquad \det A_H=1.
 \label{corpus9:xii:matriz-profunda}
\end{equation}
Su reducción módulo tres coincide con \(L_0\). La elección del
representante entero pertenece a la coordenatización: además de esa reducción, fija
la acción sobre los cocientes profundos. Para vectores fila, definimos
\begin{equation}
 \begin{aligned}
 \mathsf{Div}_H:\mathbb Z^6&\longrightarrow
       \{0,1,2\}^6\times\mathbb Z^6,\\
 x_H&\longmapsto(u_H,x_H^+),\\
 y_H&=x_HA_H,\qquad u_H=[y_H]_3,\qquad
 x_H^+=(y_H-u_H)/3.
 \end{aligned}
 \label{corpus9:xii:division-profunda}
\end{equation}
La reducción \([\,\cdot\,]_3\) se efectúa componente a componente con
representantes \(0,1,2\), también para componentes enteras negativas.

\begin{proposition}[Inversión exacta del residuo y el cociente]
\label{corpus9:xii:inversion-profunda}
La aplicación \(\mathsf{Div}_H\) es biyectiva. Su inversa es
\begin{equation}
 x_H=(u_H+3x_H^+)A_H^{-1}.
 \label{corpus9:xii:inversa-profunda}
\end{equation}
Para una trayectoria de esta coordenatización,
\(x_jA_H=u_j+3x_{j+1}\), la reconstrucción a profundidad
\(T\geq1\) satisface
\begin{equation}
 x_0=\sum_{j=0}^{T-1}3^ju_jA_H^{-(j+1)}
                 +3^Tx_TA_H^{-T}.
 \label{corpus9:xii:reconstruccion-profunda}
\end{equation}
\end{proposition}
\begin{proof}
El determinante de \eqref{corpus9:xii:matriz-profunda} es uno, por lo
que la fórmula de la adjunta da \(A_H^{-1}\in M_6(\mathbb Z)\).
La división euclídea de cada componente de \(x_HA_H\) determina un
único par \((u_H,x_H^+)\). Recíprocamente, para cualquier
\((u,q)\in\{0,1,2\}^6\times\mathbb Z^6\), el vector
\(x=(u+3q)A_H^{-1}\) es entero y cumple \(xA_H=u+3q\).
Su residuo es exactamente \(u\) y su cociente es \(q\). Esto prueba
ambas composiciones inversas. La identidad
\(x_j=(u_j+3x_{j+1})A_H^{-1}\), sustituida sucesivamente para
\(j=0,\ldots,T-1\), da
\eqref{corpus9:xii:reconstruccion-profunda}; el último término conserva
la coordenada terminal que queda después del prefijo leído.
\end{proof}

La orientación ternaria se lee posteriormente mediante
\(0\mapsto0\), \(1\mapsto+1\) y \(2\mapsto-1\), conservando
el residuo \(u_H\) y el cociente de
\eqref{corpus9:xii:division-profunda}. Los restantes campos del estado
mantienen sus propios registros de orientación, acarreo, ruta y frontera.
La proposición~\ref{prop:ch-hensel-bruto} establece, con sus pruebas de
compatibilidad y paso al límite, la torre aritmética correspondiente a
\(p=3\), \(d=6\) y \(A=A_H\). La realización de las prolongaciones
generadas utiliza además las compatibilidades APP--TRIT--TPK explicitadas
en la proposición~\ref{prop:cobertura-729-seccion}.

El bloque residual conserva también su índice
\[
 \nu(u_H)=\sum_{j=1}^{6}(u_H)_j3^{6-j}\in\{0,\ldots,728\}.
\]
Este índice se incorpora a los prefijos mediante
\(N'=729N+\nu(u_H)\). La coordenatización mixta de la
sección~\ref{act21:subsec:df-cilindros-mixtos} conserva simultáneamente
los índices ternario y decimal y los enteros de frontera
\(A(s),B(s)\). Sus recurrencias y su compatibilidad con el truncamiento
se demuestran en la proposición~\ref{act21:prop:df-rec-frontera} y el
teorema~\ref{act21:thm:df-refinamiento-generado}. Así quedan enlazadas la
recuperación de la coordenada entera y la comparación exacta de los
cilindros que expresan sus prefijos.
% END REV09_MEMORIA_PROFUNDA
